\chapter{Escopo, Motivação e Estado do Projeto}

O projeto RADOME propõe uma rede distribuída de sensoriamento eletromagnético passivo baseada em radomes geodésicos e conformais instalados em sítios elevados. Cada estação combina aberturas receptoras multifaixa, digitalização local, processamento orientado a eventos, referência temporal, calibração e rede óptica. Várias estações cooperam para estimar ângulo de chegada (AOA), diferença de tempo de chegada (TDOA), diferença de frequência de chegada (FDOA), Doppler e polarização.

O termo \emph{radome} deriva de \emph{radar dome}: um invólucro colocado sobre uma antena para protegê-la de vento, temperatura, poeira, chuva e neve, procurando preservar sua transparência eletromagnética \cite{qamar2020}. Neste projeto, porém, o domo não abriga um transmissor de radar próprio. Ele protege e integra exclusivamente aberturas receptoras. A função de radar emerge da cooperação entre estações: cada nó observa sinais diretos de emissores ou sinais de iluminadores externos refletidos por alvos, e a rede estima posição e movimento por triangulação de AOA e por multilateração/fusão de TDOA, FDOA e Doppler. Portanto, trata-se de um radar passivo e multiestático baseado em radomes receptores, e não de radares ativos distribuídos.

Essa operação exclusivamente receptora torna a estação substancialmente menos observável no domínio de emissões do que um radar ativo: não existem pulsos de iluminação, amplificador de potência ou varredura RF próprios que revelem sua posição e seu modo de operação. O sistema escuta formas de onda externas de origem e geometria conhecidas; qualquer energia atribuível à sua presença por um sensor RF externo é espalhamento secundário, isto é, eco da iluminação ambiente, e não emissão de interrogação produzida pelo radome. Essa discrição eletromagnética é uma propriedade arquitetural. Ela é distinta da seção reta radar da estrutura, cuja magnitude depende de materiais e geometria.

Muitos radomes convencionais são desenvolvidos para radares ativos, nos quais a cobertura deve preservar, nos dois sentidos, a onda transmitida e o eco recebido. Cobertura larga, varredura angular e polarização impõem co-projeto de permissividade, espessuras, núcleos e múltiplas lâminas, além de fabricação e validação eletromagnética especializadas \cite{qamar2020,ramanamurthy2016}. A arquitetura RADOME investiga outra relação custo--desempenho: as Yagis receptoras ficam expostas e orientadas pela normal de cada face, de modo que seu caminho útil não depende da transparência multifaixa do casco. Isso pode simplificar a caracterização e a substituição por módulo, mas transfere exigências para suportes externos, vedação, cabos, corrosão, vento, descargas e recalibração. Menor custo de aquisição ou de ciclo de vida é, portanto, uma hipótese a quantificar, não um resultado já demonstrado.

A rede pode explorar iluminadores de oportunidade conhecidos, incluindo infraestrutura celular, radiodifusoras, satélites e outras fontes de RF autorizadas. Na operação biestática, o caminho direto de referência é comparado ao caminho que envolve o alvo e o receptor. O projeto é conceitual: desempenho, cobertura, locais e precisão devem ser demonstrados experimentalmente e não constituem garantias operacionais.

Os objetivos são desenvolver uma arquitetura multifaixa mensurável, validar um demonstrador de três nós VHF/UHF, caracterizar os atrasos temporais e RF ponta a ponta, preservar janelas de eventos localmente e integrar a plataforma eletromagnética à infraestrutura energética, térmica, de comunicações e de sítio.

A edição independente em português brasileiro, montada a partir de \texttt{radome-pt-br.tex}, é uma das duas publicações técnicas autoritativas. A ``versão do documento 1.1'' identifica a publicação controlada, enquanto a ``revisão da arquitetura 3'' identifica a arquitetura conceitual vigente; os dois identificadores são independentes. Documentos Markdown, PDF e LaTeX autônomos anteriores são fontes históricas, salvo quando seu conteúdo for incorporado explicitamente por um capítulo vigente. Os valores quantitativos são controlados em \texttt{PARAMETERS.md}, com estado proposto, simulado, medido ou derivado, e as decisões de arquitetura são registradas em \texttt{DECISIONS.md}. Em caso de conflito, uma entrada aprovada do registro de parâmetros e este artigo prevalecem sobre o material histórico.
